% ===========================================================================
% FIGURE — the lifecycle: the eight teardown triggers and the retire/reap
% split.  Source: docs/design/l1_os_shell.md §7.6 (T0..T7) and §7.8 (the
% conservation ledger), cross-checked against kayfabe-core/src/gpu.rs
% (Proc::retire, Gpu::reap_retired) and kayfabe-isolate (two-stage retire).
% Status colours are this document's reading of the tree, not the doc's.
% ===========================================================================
\begin{tikzpicture}[x=1mm,y=1mm, font=\sffamily\scriptsize]
\def\tc#1{{\ttfamily\fontsize{5.4}{6}\selectfont #1}}

% --------------------------------------------------------- the central law --
\node[box, fill=kfcore, draw=kfcoreln, line width=1.1pt,
      minimum width=176mm, minimum height=11mm] at (0,72)
  {\textbf{The law every path obeys: RETIRE EAGER, REAP DEFERRED}\\[2pt]
   \tiny \textbf{Retire} is immediate and cheap: the \tc{Proc} stops accepting operations, its isolates refuse new
   check-outs. \;\textbf{Reap} is the heavy half — freeing host objects,\\[-0.5pt]
   \tiny killing isolate processes, recycling the GPA arena — and it runs later, at a quiesce point the adapter
   declares (\tc{Gpu::reap\_retired}).\\[-0.5pt]
   \tiny \textbf{Why:} the C artifact proved that reaping at the client-root free \emph{hangs the dying context's own
   residual polls}. Reaping never is the \#80 leak. Both were paid for on hardware.};

% ---------------------------------------------------------- the eight rows --
\def\row#1#2#3#4#5{%
  \node[box, fill=white, draw=kfportln, minimum width=26mm, minimum height=10mm,
        anchor=west, font=\sffamily\tiny]
    at (-88,#1) {\textbf{\normalsize #2}\\[1.5pt]#3};
  \node[anchor=west, align=left, text width=94mm, font=\sffamily\tiny] at (-60,#1) {#4};
  \node[anchor=west, align=left, text width=42mm, font=\sffamily\tiny] at (38,#1) {#5};
}

\node[anchor=west, font=\sffamily\bfseries\tiny, color=kfgrey] at (-88,63) {TRIGGER};
\node[anchor=west, font=\sffamily\bfseries\tiny, color=kfgrey] at (-60,63) {WHAT HAPPENS};
\node[anchor=west, font=\sffamily\bfseries\tiny, color=kfgrey] at (38,63) {STATUS — read from the tree};
\draw[kfrule] (-88,61) -- (88,61);

\row{55}{T0}{a guest frees a\\ subset and keeps\\ running}
  {The most frequent path in a real workload, and \textbf{the one with no
   backstop} — nothing dies, so nothing reclaims. Dropped host identities move to a
   \tc{pending\_release} queue on the \tc{Proc} and drain on a checked-out worker.}
  {\textbf{\color{kfstubln}Built, and lazier than designed.} The drain may fire only when
   the isolate is otherwise idle (\tc{is\_quiesced}), because no lock can exclude a verb that runs lock-free.}

\row{43}{T1}{a verb is cancelled\\ mid-flight}
  {The cancel is latched, discharged with no lock held, surfaces as \tc{EINTR} $\rightarrow$
   \tc{Interrupted} $\rightarrow$ \tc{FwdFault::Cancelled}; the chain's own unwind frees what it had
   allocated; the worker is checked back in.}
  {\textbf{\color{kfcoreln}Built and tested} (\tc{tests/tests/cancellation.rs}).}

\row{31}{T2}{a guest process\\ exits — normally,\\ killed, or killed\\ mid-verb}
  {All three are one path: the guest kernel frees the client root on fd close, so an
   \tc{RmEvent} arrives, the projection finds no boundary, and the retire is
   \textbf{graph-driven}. Every checked-out worker is cancelled; fresh handles become
   \tc{Orphans}; the \tc{Proc} lands on the retired list awaiting T3.}
  {\textbf{\color{kfcoreln}Built and tested.}}

\row{19}{T3}{the GPU goes idle —\\ THE quiesce point}
  {The guest \emph{tells} us: its driver emits \tc{UNLOADING\_GUEST\_DRIVER} (RPC fn 47)
   on an idle release and re-runs the queue handshake. That re-handshake is the quiesce
   point, and the C already ran its deferred reap there. Same doctrine as the address
   table: an observed protocol event, never an inference about idleness.}
  {\textbf{\color{kfhostln}HALF A LAW.} \tc{reap\_retired} exists and
   \tc{CoreEvent::Deferred(DeferredReap)} exists — \textbf{and nothing arms either}.
   Reap-deferred with no trigger is indistinguishable from reap-never in a run that never quiesces.}

\row{5}{T4}{the guest driver\\ restarts / the GSP\\ re-handshakes}
  {On \tc{rmmod}, guest panic or VM reset the guest sends \emph{no} \tc{Free} events, so
   the graph, every live \tc{Proc}, its isolates, arenas, host handles and routing maps
   all survive into the next driver life — which then derives a second set of components
   beside the corpses.}
  {\textbf{\color{kfhostln}NOT BUILT, and the obvious trigger does not exist.}
   Bench-measured 2026-07-26: \tc{rmmod} emits \textbf{no fn-47 at all} — the idle release
   at process exit already consumed it. \textbf{The unload has no observable signal.}}

\row{-9}{T5}{an isolate worker\\ dies}
  {Routing is settled (the component is condemned, keyed on its client set). Reclamation
   is not: a worker that died \emph{mid-chain} left allocations that are in no \tc{Orphans}
   and in no core state. The answer is to escalate to isolate death — \tc{SIGKILL} — because
   killing the process is how you reclaim state you cannot enumerate.}
  {\textbf{\color{kfstubln}Partly built.} The design calls this
   "the single most leak-prone moment"; the carrier type (\tc{VerbFailure\{err, orphans\}}) exists.}

\row{-21}{T6}{isolate garbage\\ collection}
  {\tc{SIGKILL} $\rightarrow$ the process's \emph{exit descriptor} becomes readable $\rightarrow$
   reactor $\rightarrow$ \tc{waitpid} (non-blocking) $\rightarrow$ close descriptors, tear down
   every double-mmap, release arenas, drop worker slots. The kernel already released the host
   RM objects by closing the connection.}
  {\textbf{\color{kfstubln}Designed; the OS half is unexercised} — no real isolate process
   has ever been spawned.}

\row{-33}{T7}{whole-device\\ teardown}
  {An explicit ordered \tc{shutdown()}, never a \tc{Drop}: stop accepting traps, cancel every
   worker, wait bounded, retire everything, reap unconditionally, kill every isolate, drain the
   inbox asserting every remaining event faults rather than mutates, join the threads with no
   lock held, and \textbf{assert the conservation ledger balances}. \tc{Drop} is a tripwire that
   logs loudly if \tc{shutdown()} was skipped.}
  {\textbf{\color{kfstubln}Designed; partly built.}}

% ------------------------------------------------------------- the ledger ---
\node[box, fill=kfshell!45, draw=kfshellln, minimum width=176mm, minimum height=12mm] at (0,-51)
  {\textbf{The instrument all eight are measured by — the conservation ledger} \;
   \tc{HostLedger}, \tc{kayfabe-mocks/src/lib.rs:1241}\\[2pt]
   \tiny Replays the recorded verb log and demands: every acquisition matched by exactly one release, nothing
   released that was never acquired — objects \emph{and} mappings.\\[-0.5pt]
   \tiny A test may declare a \tc{ResidueClaim} for what it knowingly leaves outstanding, and the counts are
   compared for \textbf{equality, not as a bound}: less residue than declared fails too, so a fix that closes a
   leak must delete the claim that excused it.\\[-0.5pt]
   \tiny \textbf{★ Its first census over the mean run reported zero leaks — and that was a true negative, not a
   pass.} The workload it was measuring only ever added and removed \emph{RPC} bindings, which own nothing
   host-side.\\[-0.5pt]
   \tiny Once a phase was written that actually allocates before it frees, the honest baseline was
   \textbf{24 objects, 6 mappings and 24\,576 GPA bytes leaking, linear in frees.}};

\end{tikzpicture}
